% Extension to insert after the thermal block of hibridacion.tex.
% Provenance, dependencies and scope of the checks: INFORME.md.
\section{Marked Reading and the Thermometric Conjugacy of Action}
\label{md:r27:termica}

Oriented action, the calendar, the capacity readers and the memory archive
arise from the same enriched APP--TRIT--TPK dynamics. The two sheets preserve
remainder and quotient; the tritic regime preserves orientation; transport
$U_t$ and its holonomy $\Gamma_9$ precede the reduced reading $K_{\rm ph}$.
Phase return preserves a history that continues to grow. The present
construction composes these results in a reference-and-test register, within
the previously constructed joint discrete structure of the continuum.
Its thermometric rule is specified by two operations on that register.

\subsection{Capacity, Spectral Reference and Retained Origin}
\label{md:r27:capacidad}
The capacity of a prefix indexed by $t\ge0$ and its remainder are
\begin{equation}
 C_t=\lfloor\log_{10}(9^{t+1})\rfloor,\qquad
 9^{t+1}=10^{C_t}R_t,\qquad 1\le R_t<10.
 \label{md:r27:capacidad-resto}
\end{equation}
The count is evaluated through integer powers. The inequalities
$10^{20}\le9^{21}<10^{21}$ and $10^{20}\le9^{22}<10^{21}$,
together with $10^{j-1}\le9^j<10^j$ for $1\le j\le21$, prove that
the first positive vacancy has index $21$ and capacity $20$.
In the tritic reading of capacity, residues $0,1,2$ correspond to
$+1,-1,0$, respectively. The capacities at $t=21,22,23,24$ are
$20,21,22,23$; the first subsequent return of the regime is therefore
$t=24$. The prefix contains $25$ positions and yields capacity $23$.
The distinct roles of $(20,24,25,23)$ are thereby retained.
The return of the regime retains a different nonadic phase; the phase returns
at $t=30$, with capacity $29$.

The carrier of the reader is $\mathcal D=\mathcal D_0\oplus\mathcal D_1
\oplus\mathcal D_2$, with ranks $(1,380,649)$ and grading $N=0,1,2$.
The construction preserves two gradings: $N$ distinguishes the prefix and $L$
the occupation of the exterior or symmetric space. For $m=20$, $n=24$,
let $\mathcal F_2^J$ be the degree-two exterior fixed space of two exchanged
sheets and $\mathcal B_{\le2}^J$ the symmetric fixed space through degree two.
The carrier is
$\mathbb C\Omega_{\rm ext}\oplus\mathcal F_2^J\oplus\mathcal B_{\le2}^J$.
The involution $J$ exchanges the two sheets through the graded interchange;
in the exterior sector, each of its eigenspaces has dimension $m$.
The characters of the fixed subspaces are
\begin{equation}
 \begin{aligned}
 \chi_F(z)&=\frac{(1+z)^{2m}+(1-z^2)^m}{2},\\
 \chi_B(z)&=\frac{P_n(z)^2+P_n(z^2)}{2},\\
 P_n(z)&=\prod_{j\ge1}(1-z^j)^{-n}.
 \end{aligned}
 \label{md:r27:caracteres-fock}
\end{equation}
Indeed, the projector $(I+J)/2$ averages the total character and the trace
of interchange. In the exterior space, the eigenvalues $+1,-1$ of $J$
produce $(1+z)^m(1-z)^m$; in the symmetric product, only identical pairs
contribute to the trace of interchange, producing $P_n(z^2)$.
For coefficients through degree two, it suffices to use
$P_n(z)=1+nz+n(n+3)z^2/2+O(z^3)$; the formal products are finite at
each degree. Coefficient extraction proves
\begin{equation}
 \begin{split}
 \dim\mathcal F_2^J&=m(m-1)=380,\\
 (\dim\mathcal B_0^J,\dim\mathcal B_1^J,\dim\mathcal B_2^J)
 &=(1,n,n(n+2))=(1,24,624).
 \end{split}
 \label{md:r27:grados-fock}
\end{equation}
Thus, the simultaneous decomposition $(N,L;\text{multiplicity})$ is
\[
 \begin{gathered}
 (0,0;1),(1,2;380),(2,0;1),\\
 (2,1;24),(2,2;624).
 \end{gathered}
\]
The rank $649=1+24+624=25^2+24$ retains three occupations; the external vacuum
and the vacuum of the symmetric summand remain orthogonal. The joint character
is $1+380xy^2+x^2(1+24y+624y^2)$.
Assigning $N=0,1,2$ to the cuts $24,27,30$, with capacities
$23,26,29$, specifies the reader used in this realization, while retaining $L$
as a concomitant observable. In particular,
$\operatorname{Tr}q^L=2+24q+1004q^2$ and
$\operatorname{Tr}q^N=1+380q+649q^2$ preserve distinct operations.

The spectral reference originates in
\eqref{md:antmem:cuantica:eq:5.1}--\eqref{md:antmem:cuantica:eq:5.10}.
The nonadic $8:1$ coupling and the ordered loop
$H_1=\left(\begin{smallmatrix}2&1\\1&1\end{smallmatrix}\right)$,
with the preparation of two identical pure modes and the covariance chart
transported jointly, yield
$\nu^2=1+(8/81)[\operatorname{tr}(H_1H_1^{\mathsf T})-2]=121/81$.
The sum of the squared entries of $H_1$ is seven, so that
$\nu=11/9$. Theorem~\ref{md:antmem:gauss:thm:espectro-gaussiano}
then gives $r=(\nu-1)/(\nu+1)=1/10$ and eigenvalues $(1-r)r^j$.
This established preparation is used, including its memory reduction.
The division $j=3n+b$, $0\le b<3$, preserves quotient and residue and satisfies
\begin{equation}
 (1-r)r^{3n+b}=(1-r^3)(r^3)^n\frac{r^b}{1+r+r^2}.
 \label{md:r27:factorizacion}
\end{equation}
Consequently, the quotient has ratio $q=r^3=1/1000$, while the residue retains
the distribution $\operatorname{diag}(100,10,1)/111$. Increasing the quotient
after completing the three residues preserves the carry. If the elementary
energy is $\epsilon$, the indexed energy is $\epsilon(3n+b+1/2)$: the
interblock spacing is $3\epsilon$ at the same reference temperature.

Define the positive reference operator and its intensity by
\begin{equation}
 \begin{gathered}
 \eta=\bigoplus_{n=0}^{2}r^{23+3n}I_{d_n},\qquad
 (d_0,d_1,d_2)=(1,380,649),\\
 b_*:=\operatorname{Tr}\eta=\frac{1380649}{10^{29}}.
 \end{gathered}
 \label{md:r27:eta}
\end{equation}
In particular, $b_*=1.380649\times10^{-23}$ and $0<b_*<1$.
The origin $23$ is retained in $\eta$; normalizing its state alone
removes that factor from the relative probabilities.

The action section of \eqref{md:hib:reloj} supplies
$\epsilon_a=h_a/(108t_0)$ and $\beta_{\rm clk}=\ln3/\epsilon_a$.
For the additive energy lift of this capacity reader,
\begin{equation}
 \begin{split}
 H_{\rm tot}(t)&=2(t+1)\epsilon_a,\\
 H_{\rm mem}(t)&=\frac{\epsilon_a}{\ln3}\ln R_t,\\
 H_{\rm cap}(t)&=H_{\rm tot}(t)-H_{\rm mem}(t)
 =\frac{\epsilon_a\ln10}{\ln3}C_t.
 \end{split}
 \label{md:r27:energia-capacidad}
\end{equation}
Taking logarithms in \eqref{md:r27:capacidad-resto} proves the identity.
Substitution yields $e^{-\beta_{\rm clk}H_{\rm cap}(t)}=10^{-C_t}$.
On $\mathcal D$, this gives
$H_{\rm cap}=(\epsilon_a\ln10/\ln3)(23I+3N)$ and
$\operatorname{Tr}e^{-\beta_{\rm clk}H_{\rm cap}}=b_*$.
The remainder remains recorded as an energy contribution; compensation
is a reading distinct from eliminating a factor by partial trace.

\subsection{Tensor Identity and Specification of the Readers}
\label{md:r27:lectores}
\begin{lemma}[Incremental information of the product register]
For every positive matrix $\eta$ and density matrix $\rho$, with
$\mathscr H(X)=-\operatorname{Tr}(X\ln X)$ and $0\ln0=0$,
\begin{equation}
 \begin{gathered}
 \mathscr H(\eta\otimes\rho)-\mathscr H(\eta)
 =\operatorname{Tr}(\eta)s(\rho),\\
 s(\rho)=-\operatorname{Tr}(\rho\ln\rho).
 \end{gathered}
 \label{md:r27:tensor}
\end{equation}
\end{lemma}
\begin{proof}
On the supports, the logarithm of the tensor product is
$\ln\eta\otimes I+I\otimes\ln\rho$. Trace factorization and
$\operatorname{Tr}\rho=1$ give
$\mathscr H(\eta\otimes\rho)=\mathscr H(\eta)+\operatorname{Tr}(\eta)s(\rho)$.
Continuity of $x\ln x$ at zero extends the identity to nontrivial kernels.
\end{proof}

For \eqref{md:r27:eta}, form a normalized register and its marker:
\begin{equation}
 \widehat\eta=\eta\oplus(1-b_*)|f\rangle\langle f|,
 \qquad\Omega_\rho=\widehat\eta\otimes\rho,\qquad
 P=I_{\mathcal D}\oplus0,
 \label{md:r27:registro}
\end{equation}
where $P$ acts as the identity on the test factor. The flag $f$ completes
the probability of this realization. Any additional microscopic archive
is transported as such; the flag preserves only this selection distinction.
For a fixed Hamiltonian $H$, specify
\begin{equation}
 \begin{split}
 \mathcal S_P(\rho)&=-\operatorname{Tr}
 [P\Omega_\rho(I\otimes\ln\rho)]=b_*s(\rho),\\
 \mathcal E_P(\rho)&=
 \frac{\operatorname{Tr}[P\Omega_\rho(I\otimes H)]}
 {\operatorname{Tr}(P\Omega_\rho)}=\operatorname{Tr}(\rho H).
 \end{split}
 \label{md:r27:par-lector}
\end{equation}
The first operation reads information before conditioning; the second
reads energy conditioned on the marker. Both act on the same register,
and the probability of the marker remains $b_*$. The reference remains
fixed as the test state varies.

\begin{theorem}[Thermometric conjugacy of the marked reader]
Let $H$ be a fixed self-adjoint Hamiltonian on a finite-dimensional space and
$\rho_\beta=e^{-\beta H}/Z(\beta)$, with $\beta>0$ and
$\operatorname{Var}_{\rho_\beta}(H)>0$. For the pair
\eqref{md:r27:par-lector}, the conjugate temperature satisfies
\begin{equation}
 \Theta_P:=\left(\frac{d\mathcal S_P}{d\mathcal E_P}\right)^{-1},
 \qquad b_*\Theta_P=\beta^{-1}.
 \label{md:r27:transductor}
\end{equation}
\end{theorem}
\begin{proof}
Write $E(\beta)=\operatorname{Tr}(\rho_\beta H)$.
Differentiation of the finite spectral sum gives
$Z'/Z=-E$ and $E'=-\operatorname{Var}_{\rho_\beta}(H)$.
Since $s(\rho_\beta)=\beta E+\ln Z$, one has $s'=\beta E'$.
The fixed reference permits differentiation of $\mathcal S_P=b_*s$
without an additional term, and the positive variance permits division by
$E'$. Hence $d\mathcal S_P/d\mathcal E_P=b_*\beta$, proving the claim.
\end{proof}

Normalization is part of the operational definition. If information and
energy are both read conditionally, their values are $(s,E)$ and their
differential ratio is $\beta$. If both are read before conditioning, their
values are $(b_*s,b_*E)$ and the same ratio is $\beta$. The pair
\eqref{md:r27:par-lector} yields $b_*\beta$ because it retains intensity in
the first operation and uses energy per selected test in the second.
Preserving that pair fixes its coefficient: replacing it by $cb_*$, with
$c\ne1$, changes the equation. If the reference varied,
$d(b_*s)=b_*\,ds+s\,db_*$; that protocol would have a different response
and would require retaining the second term.

\subsection{Variational Selection and Information Balance}
\begin{theorem}[Variational principle in the marked realization]
For $\Theta>0$, the functional
\begin{equation}
 \Phi_\Theta(\rho)=\mathcal E_P(\rho)-\Theta\mathcal S_P(\rho)
 =\operatorname{Tr}(\rho H)-b_*\Theta s(\rho)
 \label{md:r27:funcional}
\end{equation}
has a unique minimum over density matrices, attained at
\begin{equation}
 \rho_\Theta^*=
 \frac{e^{-H/(b_*\Theta)}}{\operatorname{Tr}e^{-H/(b_*\Theta)}}.
 \label{md:r27:gibbs}
\end{equation}
Moreover,
\begin{equation}
 \Phi_\Theta(\rho)-\Phi_\Theta(\rho_\Theta^*)
 =b_*\Theta D(\rho\Vert\rho_\Theta^*)\ge0.
 \label{md:r27:variacional}
\end{equation}
\end{theorem}
\begin{proof}
The state \eqref{md:r27:gibbs} is faithful in finite dimension.
Substituting $\ln\rho_\Theta^*=-H/(b_*\Theta)-\ln Z_\Theta$ in
$D(\rho\Vert\rho_\Theta^*)=\operatorname{Tr}\rho(\ln\rho-\ln\rho_\Theta^*)$
gives the identity. To verify positivity, diagonalize $\rho$ and the faithful
state $\sigma=\rho_\Theta^*$, with eigenvalues $r_i,s_j$ and the bistochastic
overlap matrix $p_{ij}$. Concavity of $\ln$ gives
$D(\rho\Vert\sigma)\ge\sum_i r_i\ln(r_i/t_i)$,
where $t_i=\sum_jp_{ij}s_j>0$ and $\sum_i t_i=1$.
The inequality $-\ln u\ge1-u$ shows that this sum is nonnegative.
Equality requires $r_i=t_i$ at every index and each eigenvector of $\rho$
to lie in an eigenspace of $\sigma$ with eigenvalue $t_i$;
thus it requires $\rho=\sigma$. This proves the minimum and its uniqueness.
The convention on the kernel includes singular states; faithfulness of
$\sigma$ prevents such a state from attaining equality.
\end{proof}
In this realization, the same coefficient enters both differential conjugacy
and variational equilibrium. Generalization to infinite systems separately
retains the conditions for trace, differentiability, and finite energy and
entropy; the preceding theorems have finite-dimensional domains.

\subsection{Composition with Action, Area and Gravitational Reciprocity}
\label{md:r27:composicion}
Use the previously constructed oriented sections in a common chart:
\begin{equation}
 \begin{gathered}
 \epsilon_a=\frac{h_a}{108t_0}=\frac{\hbar_a c_{\rm int}}{L_\alpha},
 \qquad q_A=\gamma a_0L_\alpha^2,\\
 \mathcal T_{I/A}=\frac{\epsilon_a}{q_A},\qquad
 G_a=\frac{c_{\rm int}^3L_\alpha^2}{\hbar_a}.
 \end{gathered}
 \label{md:r27:secciones}
\end{equation}
The longitudinal identity retains the clock of
\eqref{md:rigidez:reloj}; $q_A>0$ fixes the areal orientation used.
The functional $\gamma$ is that of \eqref{md:compos:funcional-barbero} and
retains the hexad--octad incidences and the $90/120$ channels. Its occurrence
in the area preserves that complete dependency.
Section~\ref{md:rigidez:area} constructs $a_0$ as a normalized tritic angular
character and proves $\mathcal A=q_A(N_{90}+N_{120})$ on thirty-six labelled
links. Thus $q_A$ is the area increment per unit occupation in that section,
rather than a factor inferred from a predetermined entropy.

\begin{corollary}[Compatibility of thermal, areal and gravitational readers]
With the tritic reference $\beta_{\rm clk}=\ln3/\epsilon_a$ and the same pair
\eqref{md:r27:par-lector},
\begin{equation}
 \begin{gathered}
 b_*\Theta_{{\rm clk},P}
 =\frac{\epsilon_a}{\ln3}
 =\frac{\mathcal T_{I/A}\gamma a_0L_\alpha^2}{\ln3},\\
 G_a b_*\Theta_{{\rm clk},P}\ln3=c_{\rm int}^4L_\alpha.
 \end{gathered}
 \label{md:r27:area-gravedad}
\end{equation}
For the horizon section $\tau_R=\hbar_a c_{\rm int}/(2\pi R)$,
\begin{equation}
 b_*T_{R,P}=\tau_R,\qquad
 \frac{T_{R,P}}{\Theta_{{\rm clk},P}}
 =\frac{\ln3}{2\pi}\frac{L_\alpha}{R}.
 \label{md:r27:horizonte}
\end{equation}
\end{corollary}
\begin{proof}
Applying \eqref{md:r27:transductor} to the clock gives the first identity.
Substituting $q_A$ in $\mathcal T_{I/A}q_A=\epsilon_a$ gives the second.
Multiplication by $G_a$ and \eqref{md:r27:secciones} yield
$G_a\epsilon_a=c_{\rm int}^4L_\alpha$. For the horizon, apply the same
theorem with $\beta_R=1/\tau_R$ and divide by the clock identity.
\end{proof}
Conversions between sections $i,j$ preserve
$\Theta_{i,P}/\Theta_{j,P}=\tau_i/\tau_j$. Every conversion cycle closes
by telescoping cancellation. Horizon and clock temperatures agree precisely
when their thermal energy sections agree.

In an internal chart, $b_*$ is the coefficient of the transducer defined by
\eqref{md:r27:par-lector}. The thermal unit is specified through the conjugacy
between those readings and the chosen energy section. Comparison with the
application called Boltzmann in another protocol also transports its readers,
dimensional lines and bases; the numeral alone retains its scalar scope.
In particular, this extension introduces an explicit thermometric realization
and its proofs. Its rule is incorporated as such, without retrospectively
attributing it to every earlier thermal operator or identifying its internal
unit with an external metrological unit merely by name.
